\section{Locally definable construction sequences in nowhere dense classes}\label{sec:nd}
We prove \Cref{thm:nd-def-seq}, that in nowhere dense classes, we can obtain definable construction sequences of length $O(n^{1+\eps})$ for any $\eps>0$, with $O(n^\eps)$ parameters and quantifier rank $0$.


\thmnddefseq*

\begin{proof}[Proof sketch]
Let $G\in\CC$ have $n$ vertices.
We use the almost bounded \emph{weak coloring numbers} of nowhere dense classes~\cite{NesetrilM11a,gks}, through the following consequence~\cite{drabik2026flips}.
For $v\notin U$, let $S_s(v,U)$ be the set of vertices in $U$ reachable from $v$ by a path of length at most $s$ whose internal vertices avoid $U$.
There is a total order on $V(G)$ such that $|S_{r+1}(v,U)|\le O_{\CC,r,\eps}(n^\eps)$ for every prefix $U$ of this order and every $v\notin U$.
Such an order can be computed in time $O_{\CC,r,\eps}(n^{1+\eps})$ by the algorithm of Grohe, Kreutzer, and Siebertz~\cite{gks} for weak coloring orders: every vertex of $S_{r+1}(v,U)$ is weakly $(r+1)$-reachable from $v$, since it is the smallest vertex on the witnessing path.

For each prefix $U$, consider the partial graph $H_U$ on $V(G)$ whose positive resolved pairs are the edges of $G[V(G)\setminus U]$, whose negative relation is empty, and whose parts are the atomic types in $G$ over $U$.
Thus the vertices of $U$ are singleton parts, and the remaining vertices are grouped by their neighborhoods in $U$.
Starting with $U=V(G)$, remove the last vertex $u$ of $U$ one at a time.
To pass from $H_U$ to $H_{U'}$, where $U'=U\setminus\{u\}$, first resolve positively between $\{u\}$ and each part outside $U$ adjacent to $u$ in $G$, and then merge parts with the same atomic type over $U'$.
Finally, when $U$ is empty, one negative resolve within the unique part completes the sequence.
There are at most $|E(G)|$ positive resolves and $n-1$ merges, giving length $O_{\CC,r,\eps}(n^{1+\eps})$ by the sparsity of $\CC$.

To implement the construction, first compute the atomic-type partitions over the prefixes in forward order by partition refinement: when adding $u$, split off $\{u\}$ and refine the remaining parts by adjacency to $u$.
Using adjacency lists, this takes $O(n+|E(G)|)$ time in total.
Record these splits and the resulting adjacent parts outside the prefix.
Processing the prefixes in reverse order, resolve positively against the recorded adjacent parts and undo the recorded splits by merges.
Thus the sequence can be computed in time $O_{\CC,r,\eps}(n^{1+\eps})$, including the computation of the order.

To see local definability, consider two consecutive partial graphs $H,H'$ during the transition from $H_U$ to $H_{U'}$, and put $B=\Ball^{H'}_r(v)$.
Vertices of $U'$ are still isolated, so their balls are singletons and are easy to define.
Otherwise $B\subseteq V(G)\setminus U'$. Every resolved pair of $H'$ is an edge of $G$, so every vertex of $U'$ adjacent in $G$ to a vertex of $B$ belongs to $S_{r+1}(v,U')$.
All vertices of $U'$ outside this separator are therefore nonadjacent to every vertex of $B$.
Thus, on $B$, the atomic types in $G$ over $U$ are determined by adjacency and equality tests against the parameters
\[
S_{r+1}(v,U')\cup\{u\}.
\]
This set has size $O_{\CC,r,\eps}(n^\eps)$.
Each atomic-type part is therefore quantifier-free definable, and merging parts corresponds to taking disjunctions of their definitions.
The resolved edges are the original edges between vertices outside $U$, together with edges from $u$ to the parts already resolved against it, so they too have quantifier-free definitions using these parameters.
One unary predicate naming $B$ restricts these definitions to the required domain.
The final negative resolve is immediate, since just before it the positive relation is $E(G)$ and there is only one part.
Thus the same bounds apply to $H[\Ball^{H'}_r(v)]$, including the intermediate merges and resolves.
\end{proof}
